\section{The finite Hermitian Weil action from Poisson summation}

% BEGIN G75W1
\paragraph{G75W1. The trace character and uniqueness of the test.}
Let $E$ be CM with real subfield $F$ and CM type $\Psi$ of size $s$.
Fix the native even row count $n$ and column degree $d$.
Write $x^*=\bar x^{\mathsf t}$ and $\mathrm e(t)=\exp(2\pi i t)$.
A rational Schwartz--Bruhat test on $\operatorname{Mat}_{n,d}(E)$
is the restriction of a locally constant compactly supported function
on its finite adelic completion. Indeed its level describes a function
on a finite quotient of two lattices, which agrees with its completed
quotient by the Chinese remainder theorem. We use the symmetric
rational trace pairing
\[
 [x,y]=\operatorname{Tr}_{E/\mathbb Q}\operatorname{tr}(xy^*).
\]
On finite adeles use the continuous additive character $\mathrm e_f$
of $\mathbb A_{\mathbb Q,f}$ which equals $\mathrm e(t)$ on rational
$t$. Its product with $\mathrm e(-t)$ at the real place is trivial
on $\mathbb Q$. Finite Fourier transformation with the plus sign is
\[
 \widehat f(y)=\int f(x)\,\mathrm e_f([x,y])\,dx
\]
in self-dual Haar measure. On a lattice quotient this is a finite
Fourier transform, so it preserves Schwartz--Bruhat tests and their
$L^2$ inner product by character orthogonality. The original finite
inner product uses a constant multiple of this Haar measure, as is
seen by integrating an indicator of a period coset, so the same
unitarity holds in the original normalization.

The Hermitian theta sum of $f$ at $Z=(Z_\tau)$ has, for each rational
$x$, the archimedean exponential with exponent $2\pi i$ times
\[
 \sum_{\tau,k}(u_k Z_\tau\bar u_k^t+
               u_k z_k+z'_k\bar u_k^t),\qquad u_k=\tau(x_{k,\bullet}).
\]
The variables $z_k,z'_k$ at that place are a column and a row.
At a fixed positive period and arbitrary bounded elliptic arguments
a Gaussian dominates the linear terms on the denominator lattice,
so the series converges normally. It determines $f$: restrict the
two elliptic arguments to the conjugate pair giving the real linear
character $\mathrm e([x,y])$, $y\in E_\infty^{n\times d}$.
Integrating in $y$ over the torus dual to a common denominator lattice
recovers each coefficient, which is $f(x)$ times a nonzero Gaussian.
For a group element $g$ the transformed periods and elliptic
arguments range over the same whole domain by the invertible
fractional linear action. Thus at most one test satisfies the
prescribed theta transformation for $g$.
% END G75W1

% BEGIN G75W2
\paragraph{G75W2. Three explicit transformations, including inversion.}
Let $g=\left(\begin{smallmatrix}A&B\\C&D\end{smallmatrix}\right)$
preserve $J=\left(\begin{smallmatrix}0&-I\\I&0\end{smallmatrix}\right)$
over $E$. At each embedding set $S=CZ+D$, $S'=ZC^*+D^*$,
and $K=S^{-1}C$. The required identity takes $(Z,z_k,z'_k)$ to
\[
 ((AZ+B)S^{-1},(S')^{-1}z_k,z'_kS^{-1})
\]
and multiplies the original theta function by
$\prod_\tau (\det S_\tau)^n
\mathrm e(\sum_k z'_{\tau k}K_\tau z_{\tau k})$.
We prove it for generators.

For $t_\beta=\left(\begin{smallmatrix}I&\beta\\0&I\end{smallmatrix}\right)$,
$\beta^*=\beta$, take the transformed test
\[
 f(x)\,\mathrm e_f\big(-\tfrac12
      \operatorname{Tr}_{E/\mathbb Q}\operatorname{tr}(x\beta x^*)\big).
\]
On rational $x$ the half trace is the sum over the CM type, because
the matrix trace lies in $F$. Thus the phase cancels the change of
period $Z\mapsto Z+\beta$. It preserves absolute values and the test
space. For $m_a=\operatorname{diag}(a,(a^*)^{-1})$, take
\[
 f_a(x)=\left(\prod_{\tau}\det((a_\tau^*)^{-1})^n\right)f(xa).
\]
The change $x\mapsto xa$ in the sum supplies the identity with
$Z\mapsto aZa^*$. The squared modulus of the scalar is
$|N_{E/\mathbb Q}\det a|^{-n}$. Finite Haar measure under $x\mapsto xa$
has modulus $|N_{E/\mathbb Q}\det a|^{-n}$ by the product formula.
Thus the two factors cancel in the squared norm of $f_a$.

For $J$ the transformed test is $i^{sdn}\widehat f$. To compute the
constant, the real self-dual measure for the pairing
$2\Re(u\bar v^t)$ on one complex row at one place is $2^d$ times
Lebesgue measure. The negative-sign archimedean Fourier transform
of the row exponential in G75W1 at $v$ is
\[
 {i^d\over\det Z}\,
 \mathrm e\big(-(z'-v)Z^{-1}(z-\bar v^t)\big).
\]
At $Z=iY$ with $Y$ Hermitian positive, diagonalize $Y$ unitarily.
The ordinary real Gaussian integral yields $(\det Y)^{-1}$ and
completing the square gives the displayed linear terms. This
calculation extends to independent complex $z,z'$ and all $Z$ in
the Hermitian half space by holomorphy. The positive $Y$ form an
open set in the totally real space $i\operatorname{Herm}_d$ after
multiplication by $i$, which determines the holomorphic identity.

Apply Poisson summation on the rational matrix lattice, with positive
finite and negative archimedean Fourier characters as in G75W1.
The self-dual measures give covolume one. In one copy of $E$ the
Minkowski lattice of $\mathcal O_E$ has covolume
$|\operatorname{disc}E|^{1/2}$ in the stated real self-dual measure;
the finite unit lattice has volume $|\operatorname{disc}E|^{-1/2}$
since its annihilator is the inverse different. Equivalently one
can sum the ordinary lattice Poisson identity over finitely many
support and period cosets. Multiplying the row formula at all
embeddings gives the constant $i^{sdn}\prod_\tau(\det Z_\tau)^{-n}$.
Its exponent expands into $-z'Z^{-1}z$ plus the exponent at period
$-Z^{-1}$ with elliptic arguments $Z^{-1}z,z'Z^{-1}$.
Moving the common factors across gives exactly the prescribed positive
factor for $J$ with test $i^{sdn}\widehat f$. Its scalar has modulus
one, and G75W1 gives its unitarity.
% END G75W2

% BEGIN G75W3
\paragraph{G75W3. The linear action and the orientation of its displacements.}
The transformations in G75W2 generate the rational split unitary
group. The Siegel parabolic stabilizing the first coordinate summand
is generated by the $m_a$ and $t_\beta$, by block multiplication.
A swap $J$ on a single coordinate pair is also in the generated group:
the upper unipotent supported on that pair and its conjugate by the
full $J$ are the ordinary upper and lower transvections there.
Their product with parameters $-1,1,-1$ in upper, lower, upper
order is the single-pair $J$. For a maximal isotropic subspace $W$,
a Levi change sends its intersection with the first summand to a
standard coordinate subspace $I$. In $I^\perp/I$ the image of $W$
is transverse to the first summand, hence is a graph from the
second to the first, with Hermitian graph matrix by isotropy. An
upper translation removes that graph matrix. The resulting space is
obtained from the original first summand by swapping the remaining
coordinate pairs. Applying this to the image of the first summand
under any $g$ proves generation.

The prescribed scalar composes exactly. For $g_2g_1$ the denominators
satisfy $S_{21}=S_2(g_1 Z)S_1$ and $S'_{21}=S'_1 S'_2(g_1 Z)$.
The quadratic matrix satisfies
\[
 K_{21}=K_1+S_1^{-1}K_2(g_1Z)(S'_1)^{-1}.
\]
Indeed multiplication of bottom blocks gives the first identity,
and $A_1-(g_1 Z)C_1=(S'_1)^{-1}$ by $g_1Jg_1^*=J$.
Substituting this into the product's bottom-left block gives the
last formula; applying the same block identities to the primed
denominator gives the middle one. Thus successive changes give the
required transformation for their product, including the determinant
and elliptic exponential. By uniqueness in G75W1 any two words for
$g$ give the same test. This constructs the native chosen
$\omega(g)$, with $\omega(g_2g_1)=\omega(g_2)\omega(g_1)$.
The operators are linear, preserve the finite inner product, and
the identity acts trivially.

Finally write the native Heisenberg displacement as
\[
 H(p,b)f(x)=\mathrm e_f([b,x-p/2])f(x-p).
\]
It is a Schwartz--Bruhat test. The equality
\[
 \omega(g) H(p,b)=H((p,b)g^{-1})\omega(g)
\]
uses the two blocks of the row on the right as its arguments.
For $m_a$ those blocks are $(pa^{-1},ba^*)$ by invariance of the
trace pairing. For $t_\beta$ they are $(p,b-p\beta)$: expand the
quadratic chirp at $x-p$ and use $\beta^*=\beta$ to combine the two
mixed half traces. For $J$ the substitution $x=y+p$ in the finite
Fourier integral gives
$\widehat{H(p,b)f}(v)=\mathrm e_f([p,v+b/2])\widehat f(v+b)$,
so the blocks are $(-b,p)$. Composing proves covariance for every
$g$. If $g=t_\beta m_a$ is parabolic, the formulas in G75W2 give
the factor $\prod_\tau\det D_\tau^n$, the negative Hermitian chirp
for $Ba^*$, and the argument $xa$. This also fixes the signs of
the translation law used in the zero-variable Eisenstein relation.
% END G75W3
